% 第六章 链上音乐资源发现模块设计与实现

\chapter{链上音乐资源发现与检索机制设计}


\section{模块目标与总体设计}

在 FreeFlow 中，用户可能既需要传统元数据来实现广泛的音乐发现，也需要确定单首曲目的授权访问，
本模块设计了一套检索机制，资源搜索不仅考虑歌名、艺人、专辑等传统元数据检索，还兼顾了tokenId、合约地址、交易哈希、IPFS CID 和 resourceKey 等特殊字符串识别，并考虑了常见的输入偏差。 鉴于系统的原型性质，平台未积累大规模用户数据，故本模块未将用户行为等中心化平台常见的数据分析纳入考量。

本章把链上音乐资源发现模块划分成了查询归一化、候选召回、规则评分以及结果排序四个阶段，将查询处理、候选生成以及结果组织进行解耦，以提本模块的可拓展性。


\section{检索对象与索引字段设计}

\subsection{检索对象定义}

系统使用\texttt{CreatorRelease}作为检索的对象，由于链上的资源发布需要统经过真实的签名和交易，系统只会把状态为\texttt{PUBLISHED}并且\texttt{chainId}、\texttt{musicAssetAddress}以及\texttt{tokenId}这几个字段都不为空的资源纳入到检索的范围中。

\subsection{索引字段划分}

为兼顾自然语言检索与链上身份检索，本文将检索字段划分为文本字段、身份字段和辅助排序字段三类。各字段用途如表\ref{tab:ranking_fields_rewrite}所示。

\begin{table}[htbp]
    \linespread{1.5}
    \zihao{5}
    \songti
    \centering
    \caption{链上音乐资源检索字段设计}\label{tab:ranking_fields_rewrite}
    \begin{tabular}{|p{3.35cm}|p{8.6cm}|}
        \hline
        字段                                                       & 用途                    \\ \hline
        \texttt{title}                                           & 标题匹配、标题词元重合计算与模糊相似度计算 \\ \hline
        \texttt{artistName}                                      & 艺术家名称匹配               \\ \hline
        \texttt{albumName}                                       & 专辑名称匹配                \\ \hline
        \texttt{genreLabel}                                      & 流派标签匹配                \\ \hline
        \texttt{description}                                     & 长文本辅助召回               \\ \hline
        \texttt{chainId}                                         & 链网络识别                 \\ \hline
        \texttt{musicAssetAddress}                               & 音乐资源合约定位              \\ \hline
        \texttt{platformHubAddress}                              & 平台相关合约定位              \\ \hline
        \texttt{tokenId}                                         & 链上作品唯一 Token 定位       \\ \hline
        \texttt{publishTxHash}                                   & 发布交易定位                \\ \hline
        \texttt{audioCid}、\texttt{coverCid}、\texttt{metadataCid} & IPFS 内容寻址检索           \\ \hline
        \texttt{resourceKey}                                     & 资源唯一键定位               \\ \hline
        \texttt{comments}、\texttt{purchases}                     & 行为热度估计                \\ \hline
        \texttt{publishedAt}、\texttt{updatedAt}                  & 时效性估计                 \\ \hline
    \end{tabular}
\end{table}

上述字段设计体现出两项原则。其一，标题、艺术家、专辑与流派承担主要的内容发现功能；其二，Token ID、合约地址、CID 和交易哈希承担高确定性的资源定位功能。通过对这两类字段的联合组织，系统能够在歌曲搜索与链上资产搜索两类场景下保持较稳定的表现。


\section{查询处理与候选召回}

\subsection{查询归一化}

鉴于用户输入可能会受到大小写、分隔符、复制格式以及轻微拼写偏差所带来的影响，系统在进行召回之前会首先对原始查询开展归一化处理的工作。归一化流程包括以下步骤:

\begin{enumerate}
    \item 对原始查询进行 Unicode NFKC 归一化处理，以处理全半角转换
    \item 把所有字母都转换成小写形式；
    \item 把多余的空格合并起来，并且去除掉首尾位置的空白部分；
    \item 依据非字母数字分隔符来进行词元的切分；
    \item 生成用来进行紧凑比较的\texttt{compact}文本，以消除空格、连字符等影响。
\end{enumerate}

在归一化处理的基础之上，系统还会识别几类拥有明显格式特征的链上身份，像 EVM 地址、IPFS CID、Token ID、交易哈希、\texttt{resourceKey}等。

本文当前实现里采用字符串归一化以及轻量词元切分，即根据链上身份正则表达式模版切分输入，并未实现完整的中文分词和向量语义检索。 在冷启动场景下，此模式足以覆盖典型应用场景。

\subsection{候选召回策略}

候选召回采用“字段召回 + 近期补充”两阶段的策略。这里把用户的查询记作$q$，把候选集合记作$C$。系统首先会依据归一化处理之后的查询，去执行一次数据库字段层面的召回，从而拿到候选集合$C_1$；此外，系统还会引入一批近期发布的资源，把它们当作补充集合$C_2$，然后再把它与$C_1$进行合并并且去重，最终就得到候选集合$C$。

为了控制整体复杂度和整体延迟水平，系统为候选规模设定了一个上限。在当前的实现里，主召回所拿到的候选数量，会依据查询具体的命中情况去进行动态的调整，总量会被限制在 500 条以内；补充召回所引入的近期资源不超过200条


\section{结果组织规则与评分机制}

\subsection{评分机制定义}

为便于对候选资源进行统一比较，系统采用固定权重的规则化加和方式生成结果分值，并将该实现记为\texttt{chain\_resource\_rank\_v1}。对任意候选资源$d$和查询$q$，其总分定义如公式\eqref{eq:rank_score_rewrite}所示：

\begin{equation}
    Score(d,q)=S_{identity}(d,q)+S_{text}(d,q)+S_{fuzzy}(d,q)+S_{freshness}(d)+S_{popularity}(d)
    \label{eq:rank_score_rewrite}
\end{equation}

其中，$S_{identity}$表示身份匹配分，$S_{text}$表示文本匹配分，$S_{fuzzy}$表示模糊相似度分，$S_{freshness}$表示时效性分，$S_{popularity}$表示热度分。

该评分结构主要服务于结果组织过程，可对应为三类判断信息。第一类为资源定位信息，用于判断候选资源与查询在链上身份层面的对应关系；第二类为内容相关信息，用于判断候选资源与查询在标题、艺术家和专辑等内容层面的接近程度；第三类为辅助区分信息，用于在多个相关候选之间进一步比较发布时间和用户关注程度。

\subsection{身份匹配规则}

身份匹配规则面向链上资源确定性匹配，覆盖\texttt{resourceKey}、\texttt{tokenId}、\texttt{musicAssetAddress}、\texttt{publishTxHash}以及各类 CID 字段。这些字段歧义弱、唯一性较强，在总分中的权重较高。

当前实现中，\texttt{resourceKey}、\texttt{tokenId}、音乐合约地址和 IPFS CID 的精确匹配权重分别设置为 100、80、70 和 65。其权重设置依据主要来自字段的业务作用与定位强度：\texttt{resourceKey}直接对应系统内部资源标识，具有最高区分度；\texttt{tokenId}在指定合约上下文中具有明确唯一性；合约地址能够稳定定位到资源集合或资产发行主体；CID 能够定位到链下内容对象，在链上音乐场景中同样具有较强的检索价值。

\subsection{文本匹配规则}

文本匹配规则覆盖标题、艺术家、专辑、流派和描述五类字段。系统分别考虑精确匹配、前缀匹配、包含匹配和词元重合四种信号，并通过加权方式体现不同字段与不同匹配关系的重要程度。

当前实现中的文本字段权重如表\ref{tab:text_matching_weights}所示。

\begin{table}[htbp]
    \linespread{1.5}
    \zihao{5}
    \songti
    \centering
    \caption{文本字段匹配权重设置}\label{tab:text_matching_weights}
    \begin{tabular}{|l|c|c|c|c|}
        \hline
        字段  & 精确匹配 & 前缀匹配 & 包含匹配 & 词元重合 \\ \hline
        标题  & 45   & 30   & 20   & 14   \\ \hline
        艺术家 & 34   & 24   & 16   & 12   \\ \hline
        专辑  & 24   & 16   & 10   & 8    \\ \hline
        流派  & 20   & 14   & 8    & 6    \\ \hline
        描述  & 10   & 8    & 6    & 4    \\ \hline
    \end{tabular}
\end{table}

曲目的标题通常最直接反映用户的目标作品，因此具有较高权重；艺术家和专辑承担较强的检索引导作用；流派则更多服务于主题浏览和风格发现，描述字段文本长度较长、表述方式不稳定，其作用主要体现在辅助召回和弱相关区分。

\subsection{模糊匹配规则}

为增强系统对轻微拼写错误、分隔符变化和局部字符出错的适应能力，总分中还包括模糊匹配分值。本文采用基于 Levenshtein 距离的归一化字符串相似度，对查询与候选资源的\texttt{compact}文本进行比较。设归一化后的查询紧凑文本为$q_c$，候选字段紧凑文本为$d_c$，则相似度可表示为公式\eqref{eq:lev_similarity}：

\begin{equation}
    sim_{lev}(q_c,d_c)=1-\frac{Lev(q_c,d_c)}{\max(|q_c|,|d_c|)}
    \label{eq:lev_similarity}
\end{equation}

在这里，$Lev(\cdot)$代表编辑距离，而$|\cdot|$则代表字符串长度，相似度越高，说明二者字符串越接近。

由于短字符串模糊匹配容易导致高误召回率，系统仅会在\texttt{compact}文本长度不小于 3 时启用模糊匹配；若遇短字符串，系统会提高召回对应项的阈值来提升召回精准度。

\subsection{新鲜度与热度规则}

若多个候选资源都满足了文本相关或身份相关的匹配，系统则会进一步引入新鲜度。

新鲜度分采用指数衰减形式。设资源发布时间距当前时刻的天数为$ageDays$，则新鲜度定义如公式\eqref{eq:freshness_rewrite}所示：

\begin{equation}
    freshness(d)=0.5^{ageDays/45}
    \label{eq:freshness_rewrite}
\end{equation}

此公式表示，资源每经过 45 天，它的新鲜度贡献就会衰减成原值的一半，在相同的相关性分数下，系统会优先召回最近的资源。在当前实现当中，新鲜度项的最高贡献为 8 分，以防止新鲜度的贡献超过相关程度。

热度分由评论数与购买数共同决定，其归一化形式如公式\eqref{eq:popularity_rewrite}所示：

\begin{equation}
    popularity(d)=\min\left(1,\frac{\ln(1+comments+3\times purchases)}{\ln(31)}\right)
    \label{eq:popularity_rewrite}
\end{equation}

在当前实现当中，热度项的最高贡献为 7 分。
用户的购买行为被视为强烈的正面信号，因此比评论数目具有更高的权重；对数归一化处理能削弱极端高活跃曲目的对整体相关性的影响。


\section{结果说明与处理流程}

\subsection{结果说明机制}

为了便于分析设计，系统会记录检索流程中的分数，如命中的字段类型、匹配关系、各项特征分值以及它们对总分的贡献等。 例如，某一曲目因标题精确匹配获得了 45 分、艺术家前缀匹配获得了 24 分、新鲜度获得了 6.3 分、热度获得了 4.8 分等。该机制有助于误召回、定位参数等问题分析。

\subsection{处理流程}

链上音乐资源发现流程如图\ref{fig:search_ranking_pipeline}所示。

\begin{figure}[H]
    \centering
    \includegraphics[width=\textwidth]{figures/out/search_ranking_pipeline}
    \caption{链上音乐资源发现处理流程图}
    \label{fig:search_ranking_pipeline}
\end{figure}

对应的处理过程如算法\ref{algo:chain_resource_rank_rewrite}所示。

\begin{algorithm}[H]
    \caption{链上音乐资源发现处理流程}
    \label{algo:chain_resource_rank_rewrite}
    \small
    \begin{algorithmic}[1]
    \Require 用户查询$q$，返回数量$k$
    \Ensure 排序后的链上音乐资源列表
    \State 对$q$进行归一化处理，得到词元集合、\texttt{compact}文本和特殊字段识别结果
    \State 基于标题、艺术家、专辑、流派、描述及链上身份字段执行数据库召回，得到候选集合$C_1$
    \State 选取近期发布资源构成补充集合$C_2$
    \State 合并去重，得到最终候选集合$C=C_1\cup C_2$
    \For{每个候选资源$d \in C$}
        \State 分别计算身份匹配分、文本匹配分、模糊相似度分、新鲜度分和热度分
        \State 汇总总分$Score(d,q)$并记录结果解释信息
    \EndFor
    \State 过滤未命中有效相关特征的候选资源
    \State 按总分降序排序；若总分相同，则按发布时间降序排序
    \State 返回前$k$个结果
    \end{algorithmic}
\end{algorithm}


\section{复杂度分析}

设候选资源的数量为$N$，而平均参与比较的文本长度是$L$。

文本匹配分数的计算依赖字段的比较，时间复杂度可以近似地表达为$O(NL)$。模糊相似度使用 Levenshtein 距离计算，在单次比较当中的时间复杂度近似为$O(L^2)$，因此作用在整个候选集合之上时，总体的复杂度近似为$O(NL^2)$。最终排序这一步骤，运用了依靠总分的比较排序方式，它的复杂度为$O(N\log N)$。

\section{检索效果评估设计}

为了验证上述冷启动搜索排序模块，本文基于公开的 MusicBrainz sample 2026-04-01 数据库构建了一个离线评估集\cite{musicbrainzDatabase,musicbrainzDataDownload}。该评估集先从样本库里选取 48 个发行簇作为核心相关样本，再加入来自整库抽样的干扰项，最终形成了 6000 条候选资源，其中核心相关样本为 815 条，干扰项为 5185 条。

查询类型具体涵盖了 162 条查询，覆盖标题精确匹配、艺人查询、专辑查询、多词组合、别名检索、拼写扰动、紧凑输入、全角输入，以及 Token ID、合约地址、IPFS CID、\texttt{resourceKey}和交易哈希等链上身份检索，能有效测试自然语言查询和链上匹配特殊信息的检索效果。

在相关性标注方面，多结果查询采用三级相关性标注，并依据查询意图的匹配程度以及资源代表性标记为 3、2、1 三个等级；高确定性的链上身份查询则采用较浅的 spot-check 方式。当前设计下，人工深度标注查询共有 148 条，spot-check 查询共有 14 条，预计人工判断量约为 1550 条。设计能够在控制标注成本的同时，把多结果场景与精确定位场景都覆盖进去。

需要说明的是，尽管元数据来自真实数据，但链上身份字段、热度值以及部分时效性信号因链上测试成本不可控、无用户规模等问题，由模拟程序随机产生，因此它依然不能替代真实线上行为日志。

本文从排序质量、召回有效性，还有响应效率三个方面评估检索质量。其中，Precision@10 计算前 10 个结果中相关资源的比例，用于衡量搜索结果第一页的准确性；MRR 则衡量了首个高相关结果出现位置；nDCG@10 指标，其使用所参考的是分级相关性条件下的经典累积增益评价方法\cite{Jarvelin2002CG}，它用来衡量在该条件下前 10 个结果的排序质量；零结果率反映了检索请求里无返回结果的比例；平均响应时间则用来衡量检索链路的效率表。

实验包含字段包含匹配基线、仅文本评分以及完整\texttt{chain\_resource\_rank\_v1}等多组方法。本文在结果章节中主要保留三组代表性方法：数据库字段包含匹配基线、仅运用文本评分的方法，以及完整排序算法。分层对比，能够观察出从基础字段匹配过渡到完整规则排序之后，整体检索效果和响应开销的变化。有关具体的实验过程、数据分布以及结果分析，可参见第 7 章当中的表\ref{tab:ranking_query_categories}、表\ref{tab:ranking_metrics_actual}和图\ref{fig:ranking_latency_curve}。

